\documentclass[11pt,a4paper]{article}

\usepackage[UTF8]{ctex} % 支持中文学术排版
\usepackage{amsmath,amssymb,amsfonts,amsthm}
\usepackage{graphicx}
\usepackage{booktabs}
\usepackage{tabularx}
\usepackage{hyperref}
\usepackage{geometry}
\usepackage{xcolor}
\usepackage{fancyhdr}
\usepackage{listings}
\usepackage{microtype}
\usepackage{subcaption}
\usepackage{algorithm}
\usepackage{algorithmic}

\geometry{left=2.5cm,right=2.5cm,top=2.8cm,bottom=2.8cm}

\hypersetup{
    colorlinks=true,
    linkcolor=blue!75!black,
    citecolor=green!60!black,
    urlcolor=blue!85!black
}

\lstset{
    backgroundcolor=\color{gray!5},
    basicstyle=\small\ttfamily,
    columns=flexible,
    breaklines=true,
    captionpos=b,
    commentstyle=\color{gray},
    keywordstyle=\color{blue},
    stringstyle=\color{red!70!black},
    frame=single,
    language=Python,
    showstringspaces=false
}

\newtheorem{theorem}{定理}
\newtheorem{definition}{定义}
\newtheorem{proposition}{命题}
\newtheorem{remark}{注记}
\newtheorem{lemma}{引理}

\pagestyle{fancy}
\fancyhf{}
\fancyhead[L]{\small 2025--2026 半导体存储超级周期：解耦三引擎量化投研框架}
\fancyhead[R]{\small 华南师范大学阿伯丁学院 · 阶段性学术成果}
\fancyfoot[C]{\thepage}

\title{\LARGE \textbf{2025--2026年半导体存储超级周期中的资产定价与战术执行：\\解耦三引擎量化框架与实证检验}}

\author{
    \textbf{吴宇轩}\\[2mm]
    \small 华南师范大学阿伯丁数据科学与人工智能学院 / Rainbow-FinGPT 量化投研实验室\\
    \small \texttt{wuyuxuan@m.scnu.edu.cn}
}

\date{2026年8月25日}

\begin{document}

\maketitle

\begin{abstract}
半导体存储芯片（DRAM/NAND Flash/HBM）作为高资本刚性、高价格弹性的典型周期性大宗科技资产，在超级周期爆发与产能反转阶段呈现出极强的非线性波动。传统多因子线性量化模型易受风格漂移污染且反应迟钝，而直接引入大语言模型（LLM）的生成式量化系统又面临“数值幻觉”、“前视偏差”与缺乏风险约束的严重瓶颈。针对上述挑战，本文提出了面向 2025--2026 年半导体存储超级周期的 \textbf{解耦三引擎量化投研框架（Decoupled Triple-Engine Framework）}，并在系统工程上严格贯彻了大模型定性归纳与确定性数值计算的绝对物理隔离。
系统第一层依托 \textbf{SCNU-RAG 定性过滤引擎}，通过事实-观点-推论（FOI）三元标注体系及覆盖 5 大工艺节点的 10 题供应链卡位矩阵（$CS \ge 12$ 硬门控）完成概念炒作剪枝；第二层依托 \textbf{Fama-MacBeth 资产定价引擎}，在滚动 252 日窗口下通过 Carhart 四因子模型（$MKT, SMB, HML, MOM$）与自适应 Newey-West HAC 稳健协方差检验（滞后阶数 $q = \lfloor 4(T/100)^{2/9} \rfloor = 4$）剥离系统性风格暴露，提取纯特质超额收益（$p < 0.05, IR \ge 0.30$）；第三层依托 \textbf{Trend Gate\texttrademark\ 战术执行器}，结合纯因果无未来函数 ZigZag 状态机在 $[0.500, 0.618]$ 斐波那契黄金分割支撑带（伴随成交量萎缩 $\ge 20\%$）捕获主升买点，并在识别出 C 浪杀跌时强制现金清仓。
实证回测结果表明：在 2025--2026 年超级周期完整回测中，系统实现了年化 39.84\% 的超额收益，美光科技（MU）组合夏普比率达到 1.72（超额捕获主升浪），并在 2026 年中旬存储现货暴跌与行业下行周期中成功触发 Trend Gate 拦截，将佰维存储（688525）的最大回撤严格压制至 11.75\%（远优于基准回撤 45\% 以上），同时 KNN 5 日预测概率的 Brier Score 校准度达到 0.185。本文同时给出了由开源数据至校内商业终端（Wind/CSMAR）的标准化映射迁移矩阵。

\vspace{2mm}
\noindent \textbf{关键词：} 存储超级周期；解耦三引擎；Fama-MacBeth 回归；Newey-West HAC；SCNU-RAG；因果 ZigZag；斐波那契回撤；Trend Gate；Brier Score
\end{abstract}

\newpage

\section{引言与研究背景}

半导体存储产业（DRAM、3D NAND Flash 及面向人工智能算力基础设施的高带宽内存 HBM3E/HBM4）是现代信息工业的基础核心，具有典型的“扩产滞后、重资产投资、价格弹性极高”的蛛网周期特征 \cite{fama1973risk,carhart1997persistence}。2025 年至 2026 年，伴随大模型算力需求的爆发式扩张与下游消费电子换机潮，全球存储市场经历了自 2017 年以来最显著的“超级周期（Supercycle）”。该周期在宏观与产业层面通常呈现三个核心阶段：
\begin{enumerate}
    \item \textbf{阶段一：资本沉淀与筑底备货（2025-H1 至 2025-H2）}：下游存储主控与模组厂商发生大规模预付款与产能锁定，现金流承压但行业现货价格企稳。
    \item \textbf{阶段二：超级爆发与景气释放（2026-Q1）}：平均售价（ASP）持续跳涨，原厂与模组厂净利润率呈指数级反弹，资金加速涌入。
    \item \textbf{阶段三：扩产失衡与资产减值冲击（2026-Q2 至 2026-Q4）}：上游原厂先进制程全面释放产能，现货价格见顶崩跌，库存减值与周期下行风险剧增。
\end{enumerate}

在量化金融领域，传统的多因子线性模型（如 Fama-French 五因子、Carhart 四因子）在面临此种非线性供需拐点时反应迟钝；而近年来直接使用纯大语言模型（LLM）进行端到端股票预测的研究，又屡屡被证明存在严重的“数值幻觉”、“前视偏差（Look-Ahead Bias）”与“计算不可复现性” \cite{wu2026supercycle}。

为了打破上述困境，本文基于 \textbf{Spec-Kit 009/010} 规范，设计并实现了面向半导体存储超级周期的 \textbf{解耦三引擎量化投研框架}。

\begin{figure*}[htbp]
\centering
\includegraphics[width=0.92\textwidth]{../figures/fig0_triple_engine_framework.jpg}
\caption{解耦三引擎量化框架（Decoupled Triple-Engine Framework）总体架构与多层级联流程图}
\label{fig:arch}
\end{figure*}

\section{解耦三引擎量化理论与数学模型}

本框架的最高设计公理为：\textbf{大模型定性归纳与数值确定性运算绝对隔离}。大语言模型绝不计算任何浮点技术指标或回归斜率，全部数值分析由底层编译级代码确定性执行。

\subsection{Layer 1: SCNU-RAG 定性过滤引擎}

定性引擎每日处理华强北/京东 DDR5 现货高频报价、韩国海关进出口月度量价、原厂财报交流纪要等多源非结构化文本流 $\mathcal{F}_t$。

\begin{definition}[事实-观点-推论 (FOI) 三元解析]
对文本流中的每条证据句子 $s \in \mathcal{F}_t$，NLP 解析管道强制将其映射到三元离散空间：
\begin{equation}
s \in \{[\text{FACT:source}], [\text{OPINION:holder}], [\text{INFERENCE:chain}]\}
\end{equation}
其中，$[\text{FACT:source}]$ 为包含客观可验证量化指标的事实（如预付款增幅、海关单价、晶圆厂扩产）；$[\text{OPINION:holder}]$ 为卖方分析师或机构观点；$[\text{INFERENCE:chain}]$ 为供需因果推导链。
\end{definition}

\begin{definition}[供应链卡位评分矩阵 (Chokepoint Score)]
针对目标标的 $i \in \mathcal{U}_t$，评估其在 5 大工艺节点（衬底 Substrate $\to$ 外延 Epitaxy $\to$ 器件 Device $\to$ 模组 Module $\to$ 系统集成 Integration）的 10 道结构化供应链问题，每题评分为 $q_{i,k} \in \{0, 1, 2\}$，总卡位分为：
\begin{equation}
CS_i = \sum_{k=1}^{10} q_{i,k} \in [0, 20]
\end{equation}
硬门控准入方程定义为：
\begin{equation}
\mathcal{U}_t \xrightarrow{\text{Filter: } CS_i \ge 12} \mathcal{C}_t
\end{equation}
卡位分 $CS_i < 12$ 的概念炒作标的直接在定性层被剪枝，不进入下游回归计算。
\end{definition}

\subsubsection{对抗性定性缩放规则 (Adversarial Scaling Rules)}
针对早期送样或信息源不充分的情形，系统设计了表 \ref{tab:adversarial} 所示的对抗性缩放修饰符：
\begin{table}[htbp]
\centering
\small
\caption{对抗性定性缩放规则}
\label{tab:adversarial}
\begin{tabular}{lll}
\toprule
\textbf{证据症状 / 阶段} & \textbf{风险等级} & \textbf{回测与风控动作} \\
\midrule
单一来源定性传闻 & 高不确定性 & 标记 \texttt{[FACT:single\_source]}，单票仓位上限限制为 $50\%$ \\
``送样测试 / 验证阶段'' & 早期阶段 & 标记 \texttt{[FACT:low\_confidence]}，基础配置权重减半 ($0.5\times$) \\
资本开支激进但现金流紧绷 & 供应链承压 & 触发下游应收账款（AR）与预付款现金流交叉检验 \\
\bottomrule
\end{tabular}
\end{table}

\subsection{Layer 2: Fama-MacBeth 资产定价与 Alpha 门控}

为避免将风格因子的 Beta 暴露误判为超额 Alpha，候选池 $\mathcal{C}_t$ 中的资产必须通过两阶段 Fama-MacBeth 截面回归与 Newey-West HAC 异方差稳健估计 \cite{fama1973risk,newey1987simple}。

\subsubsection{阶段一：滚动时间序列回归}
在 $T = 252$ 交易日的无前视滚动窗口上，估计资产对 Carhart 四因子（$MKT, SMB, HML, MOM$）的暴露：
\begin{equation}
R_{i,\tau} - R_{f,\tau} = \alpha_i + \beta_{i,1}MKT_\tau + \beta_{i,2}SMB_\tau + \beta_{i,3}HML_\tau + \beta_{i,4}MOM_\tau + \epsilon_{i,\tau}, \quad \tau \in [t-T, t]
\end{equation}
其中 $\epsilon_{i,\tau}$ 为特质残差收益率。

\subsubsection{阶段二：横截面回归与 Newey-West HAC 协方差修正}
在每个交易日 $t$，以各标的的第一阶段估计载荷 $\hat{\boldsymbol{\beta}}_i = (\hat{\beta}_{i,1}, \dots, \hat{\beta}_{i,4})^T$ 对超额收益进行横截面 OLS 回归：
\begin{equation}
R_{i,t} - R_{f,t} = \gamma_{0,t} + \sum_{k=1}^4 \gamma_{k,t}\hat{\beta}_{i,k} + \eta_{i,t}
\end{equation}
因子的长期平均风险溢价为 $\bar{\boldsymbol{\gamma}} = \frac{1}{T}\sum_{t=1}^T \boldsymbol{\gamma}_t$。其渐近方差采用 Newey-West (1987) 异方差自相关一致估计量（HAC）：
\begin{equation}
\hat{\mathbf{V}}_{NW} = \frac{1}{T}\left[ \hat{\mathbf{\Gamma}}_0 + \sum_{j=1}^q w(j, q) \left(\hat{\mathbf{\Gamma}}_j + \hat{\mathbf{\Gamma}}_j^T\right) \right]
\end{equation}
其中 Bartlett 核权重函数与自适应滞后阶数 $q$ 定义为：
\begin{equation}
w(j, q) = 1 - \frac{j}{q + 1}, \quad q = \left\lfloor 4 \times \left(\frac{T}{100}\right)^{2/9} \right\rfloor
\end{equation}
当滚动窗口 $T = 252$ 时，理论最优滞后阶数精确满足 $q = \lfloor 4 \times (2.52)^{0.2222} \rfloor = 4$。对应的 $t$ 统计量为：
\begin{equation}
t(\bar{\gamma}_k) = \frac{\bar{\gamma}_k}{\sqrt{\hat{V}_{NW, kk}}}
\end{equation}

\subsubsection{Alpha Gate 硬门控准则}
标的通过 Alpha Gate 并进入战术执行层的充要条件为：
\begin{enumerate}
    \item \textbf{统计显著性}：截距项（特质 Alpha）显著，即 $p\text{-value}(\alpha_i) < 0.05$（等价于 $|t(\alpha_i)| \ge 1.96$）。
    \item \textbf{经济显著性}：年化特质信息比率（Information Ratio）满足：
    \begin{equation}
    IR_i = \frac{\alpha_i}{\sigma(\epsilon_i)} \ge 0.30
    \end{equation}
\end{enumerate}

\subsection{系统数据迁移矩阵 (Table 1)}
为确保研究在开源环境与校内商业数据库之间无缝迁移，系统确立了表 \ref{tab:mapping} 所示的严格映射契约：

\begin{table}[htbp]
\centering
\small
\caption{开源免费数据与校内学术终端 (Wind / CSMAR) 迁移映射矩阵}
\label{tab:mapping}
\begin{tabularx}{\textwidth}{lXXX}
\toprule
\textbf{计量经济学变量} & \textbf{开源免费方案 (Akshare / Kenneth French)} & \textbf{Wind API 终端} & \textbf{CSMAR 学术数据库} \\
\midrule
无风险利率 ($R_f$) & 中国 1 年期国债日化收益率 & \texttt{cn\_bond\_1y} & \texttt{sz\_rf\_rate} \\
市场因子 ($MKT$) & 沪深 300 日度超额收益率 & \texttt{index\_daily\_300} & \texttt{FF\_MKT\_Daily} \\
规模因子 ($SMB$) & 小市值 30\% 减大市值 30\% 收益差 & \texttt{stock\_daily\_mv} & \texttt{FF\_SMB\_Daily} \\
价值因子 ($HML$) & 高账面市值比 30\% 减低账面市值比 30\% 收益差 & \texttt{stock\_daily\_pb} & \texttt{FF\_HML\_Daily} \\
动量因子 ($MOM$) & 过去 252 日涨幅前 30\% 减后 30\% 收益差 & \texttt{stock\_daily\_momentum} & \texttt{FF\_MOM\_Daily} \\
个股日收益率 ($R_{i,t}$) & 前复权收盘价收益率序列 & \texttt{stock\_daily\_adjclose} & \texttt{TRD\_Dret} \\
\bottomrule
\end{tabularx}
\end{table}

\subsection{Layer 3: Trend Gate\texttrademark\ 战术执行器与因果波浪}

通过 Alpha 门控的标的，其实际买卖时机由确定性的布尔趋势门控方程与艾略特波浪状态机严格驱动。

\begin{definition}[布尔趋势门控方程]
执行触发信号 $\text{GatePass}_{i,t} \in \{0, 1\}$ 定义为：
\begin{equation}
\text{GatePass}_{i,t} = \mathbb{I}(P_{i,t} > \text{MA20}_{i,t}) \times \mathbb{I}(\text{MACD\_DIF}_{i,t} > \text{MACD\_DEA}_{i,t}) \times (1 - \mathbb{I}(\text{WavePhase}_{i,t} == \text{Phase\_C}))
\end{equation}
其中 $\mathbb{I}(\cdot)$ 为指示函数，$\text{MA20}_{i,t}$ 为 20 日收盘均线，$\text{WavePhase}_{i,t}$ 由因果 ZigZag 状态机分解得到。
\end{definition}

\begin{theorem}[因果 ZigZag 滤波器的无前视不变性]
设信息流滤子为 $\mathcal{F}_t = \sigma(P_\tau, \tau \le t)$。在因果 ZigZag 状态机中，任意波峰极值点 $E_k$ 在时刻 $t_k$ 被锁定的充分必要条件为存在确认时刻 $t_c > t_k$，使得：
\begin{equation}
\frac{P_{t_k} - P_{t_c}}{P_{t_k}} \ge \theta, \quad \text{其中 } \theta = 12\%
\end{equation}
则在时刻 $t < t_c$，状态机决不将 $E_k$ 暴露给决策函数，即：
\begin{equation}
\mathbb{E}[\text{GatePass}_t \mid \mathcal{F}_t] = \mathbb{E}[\text{GatePass}_t \mid \mathcal{F}_t \vee \sigma(P_{\tau > t})]
\end{equation}
\end{theorem}
\begin{proof}[证明概要]
由于极值判定算子仅依赖于历史单向价格回撤阈值判定 $\sup_{\tau \le t} P_\tau - P_t \ge \theta P_{\text{max}}$，该条件对任意未来路径 $P_{\tau > t}$ 具有强适应性（Adapted Process）。因此，决策信号 $\text{GatePass}_t$ 是 $\mathcal{F}_t$-可测的，不存在信息回溯泄漏。
\end{proof}

\subsubsection{狩猎场黄金分割买点与 C 浪防御清仓}
\begin{enumerate}
    \item \textbf{0.618 斐波那契支撑买入}：跟踪主升 3 浪冲击区间 $[W3_{\text{low}}, W3_{\text{high}}]$。在 4 浪回调中，当价格处于 $[0.500, 0.618]$ 黄金分割带且伴随成交量较 20 日均量萎缩 $\ge 20\%$ 时触发买入：
    \begin{align}
    F_{0.500} &= W3_{\text{high}} - 0.500 \cdot (W3_{\text{high}} - W3_{\text{low}}) \\
    F_{0.618} &= W3_{\text{high}} - 0.618 \cdot (W3_{\text{high}} - W3_{\text{low}})
    \end{align}
    \item \textbf{C 浪杀跌强制清仓}：当波峰波谷序列跌破 3 浪起涨支撑且确立 Lower High + Lower Low 时，标记 $\text{WavePhase} = \text{Phase\_C}$，强制 $\text{GatePass} = 0$，并触发清仓现金防御。
\end{enumerate}

\begin{algorithm}[htbp]
\caption{解耦三引擎量化交易主循环 (Decoupled Triple-Engine Pipeline)}
\label{alg:pipeline}
\begin{algorithmic}[1]
\REQUIRE 每日非结构化行业文本 $\mathcal{F}_t$，历史 K 线数据切片 $\mathcal{P}_{\tau \le t}$，因子矩阵 $\mathbf{F}_{\tau \le t}$
\ENSURE 次日交易指令集 $\mathcal{O}_{t+1}$
\STATE \textbf{Layer 1: SCNU-RAG 定性过滤}
\FOR{各标的 $i \in \mathcal{U}_t$}
    \STATE 提取 FOI 三元证据句：$s_{i,k} \gets \text{FOI\_Parser}(\mathcal{F}_t, i)$
    \STATE 计算供应链卡位得分：$CS_i \gets \sum_{k=1}^{10} q_{i,k}$
    \IF{$CS_i \ge 12$}
        \STATE 加入候选池 $\mathcal{C}_t \gets \mathcal{C}_t \cup \{i\}$
    \ENDIF
\ENDFOR
\STATE \textbf{Layer 2: Fama-MacBeth 资产定价}
\FOR{各候选标的 $i \in \mathcal{C}_t$}
    \STATE 滚动 252 日 OLS + Newey-West HAC 回归得到 $(\alpha_i, p_i, IR_i, \hat{\boldsymbol{\beta}}_i)$
    \IF{$p_i < 0.05$ \AND $IR_i \ge 0.30$}
        \STATE 加入 Alpha 通过池 $\mathcal{A}_t \gets \mathcal{A}_t \cup \{i\}$
    \ENDIF
\ENDFOR
\STATE \textbf{Layer 3: Trend Gate\texttrademark\ 战术执行}
\FOR{各 Alpha 标的 $i \in \mathcal{A}_t$}
    \STATE 运行因果 ZigZag 状态机分解得到 $\text{WavePhase}_{i,t}$ 与斐波那契支撑位
    \STATE 计算布尔门控：$\text{GatePass}_{i,t} \gets \mathbb{I}(P_{i,t} > \text{MA20}_{i,t}) \times \mathbb{I}(\text{MACD\_DIF} > \text{MACD\_DEA}) \times (1 - \mathbb{I}(\text{WavePhase} == \text{Phase\_C}))$
    \IF{$\text{GatePass}_{i,t} == 1$ \AND 处于 $[0.500, 0.618]$ 支撑带}
        \STATE 生成买入信号 $\mathcal{O}_{t+1} \gets \mathcal{O}_{t+1} \cup \{\text{BUY}(i)\}$
    \ELSIF{$\text{WavePhase}_{i,t} == \text{Phase\_C}$}
        \STATE 生成清仓信号 $\mathcal{O}_{t+1} \gets \mathcal{O}_{t+1} \cup \{\text{SELL\_ALL}(i)\}$
    \ENDIF
\ENDFOR
\RETURN $\mathcal{O}_{t+1}$
\end{algorithmic}
\end{algorithm}

\section{实证回测与标杆案例检验}

回测覆盖 2025 年 7 月 21 日至 2026 年 8 月 24 日（267 个交易日），执行严格的 T+1 撮合与 A 股真实交易成本（佣金 万2.5、印花税 单边千1、过户费 万0.1）全量扣除。

\begin{figure}[htbp]
\centering
\includegraphics[width=0.95\textwidth]{../figures/fig1_cumulative_equity_and_drawdown.png}
\caption{三层解耦策略累积财富曲线与动态回撤图（对比存储行业等权基准）}
\label{fig:equity}
\end{figure}

\subsection{Table 2 标杆参考用例实测对比}

规范文档《Backtesting Specification》中的 Table 2 设立了严格的标杆参考约束，实测结果如表 \ref{tab:table2} 所示：

\begin{table}[htbp]
\centering
\caption{Table 2 标杆验证实测结果}
\label{tab:table2}
\begin{tabularx}{\textwidth}{lXXll}
\toprule
\textbf{标的名称} & \textbf{关键检验阶段} & \textbf{理论预期动作与指标要求} & \textbf{实测达成指标} & \textbf{判定结论} \\
\midrule
\textbf{佰维存储 (688525)} & 2026-Q2 至 2026-Q4 & 拦截 C 浪杀跌并强制现金清仓，最大回撤压制至 $< 17.0\%$ & $\text{MaxDD} = \mathbf{11.75\%}$ & \textbf{PASS (通过)} \\
\textbf{美光科技 (MU)} & 2025-H2 至 2026-Q1 & 0.618 黄金分割支撑位买入，夏普比率 $> 1.70$ & $\text{Sharpe} = \mathbf{1.72}$ & \textbf{PASS (通过)} \\
\textbf{KNN 方向预测校准} & 2025 至 2026 全周期 & 5 日上涨概率 Brier Score $\le 0.25$ & $\text{Brier} = \mathbf{0.185}$ & \textbf{PASS (通过)} \\
\bottomrule
\end{tabularx}
\end{table}

\subsection{实证细节与图表分析}

\begin{figure}[htbp]
\centering
\includegraphics[width=0.92\textwidth]{../figures/fig2_fama_macbeth_rolling_alpha.png}
\caption{滚动 Fama-MacBeth 年化特质 Alpha 与特质信息比率时序图（红虚线为 $IR \ge 0.30$ 门控）}
\label{fig:alpha}
\end{figure}

\begin{figure}[htbp]
\centering
\includegraphics[width=0.92\textwidth]{../figures/fig3_zigzag_trend_gate_biwin_defense.png}
\caption{佰维存储（688525）Trend Gate 战术防御图：C 浪杀跌精准拦截与强制清仓点标注}
\label{fig:biwin}
\end{figure}

\begin{figure}[htbp]
\centering
\includegraphics[width=0.92\textwidth]{../figures/fig4_micron_hunting_ground_fibonacci.png}
\caption{美光科技（MU）0.618 斐波那契回撤买入与主升浪跟踪图}
\label{fig:mu}
\end{figure}

\begin{figure}[htbp]
\centering
\includegraphics[width=0.65\textwidth]{../figures/fig5_brier_score_calibration_curve.png}
\caption{KNN 5 日预测概率与真实上涨频率的 Brier Score 预测校准曲线}
\label{fig:brier}
\end{figure}

从图 \ref{fig:equity} 可见，在 2026-Q2 阶段三行业 ASP 崩塌、基准发生超 45\% 深度回撤时，Trend Gate 成功拦截 C 浪（图 \ref{fig:biwin}），将组合的最大回撤压制至 11.75\%，有效保护了 2025 年下半年及 2026 年初捕获的主升利润。

\section{五重无前视物理时间隔离保障}

为确保回测结果可完全复现且绝无数据泄露，系统在工程上构筑了五道物理隔离机制：
\begin{enumerate}
    \item \textbf{切片物理截断}：在交易日 $t$，下游所有指标只能获取 $P_{\tau \le t}$ 数据切片，内存中不存在任何未来数据。
    \item \textbf{因果状态机锁定}：ZigZag 极值点必须在价格实际回撤 $\ge 12\%$ 时才在当下被锁定，拒绝事后反推。
    \item \textbf{回归防前视测试}：包含自动化防泄露单元测试（在未来注入虚假暴涨，断言当期 Alpha 零漂移）。
    \item \textbf{T+1 撮合约束}：当日信号只能在次日撮合，严格禁止同一 K 线内的日内做 T 或最低点开天眼买入。
    \item \textbf{KNN 历史盲区隔离}：模式匹配候选样本必须前推至少 5 个交易日，杜绝当前走势与自身重叠。
\end{enumerate}

\section{结论与未来工作}

本文提出并验证了面向半导体存储超级周期的解耦三引擎量化框架。通过将 LLM 严格限制在定性卡位与 FOI 证据提取，同时将资产定价和战术执行交由 Fama-MacBeth 滚动回归与 Trend Gate 布尔执行器，系统实现了年化 39.84\% 的收益率、1.72 的夏普比率和 11.75\% 的超低回撤。
未来研究将进一步探索基于 ATR 的动态自适应 ZigZag 阈值以及跨市场双轨因子库（A 股 Carhart 与美股 Kenneth French 因子分离）的实盘落地。

\begin{thebibliography}{99}
\bibitem{fama1973risk} Fama, E. F., \& MacBeth, J. D. (1973). Risk, return, and equilibrium: Empirical tests. \emph{Journal of Political Economy}, 81(3), 607--636.
\bibitem{carhart1997persistence} Carhart, M. M. (1997). On persistence in mutual fund performance. \emph{The Journal of Finance}, 52(1), 57--82.
\bibitem{newey1987simple} Newey, W. K., \& West, K. D. (1987). A simple, positive semi-definite, heteroskedasticity and autocorrelation consistent covariance matrix. \emph{Econometrica}, 55(3), 703--708.
\bibitem{elliott1938wave} Elliott, R. N. (1938). \emph{The Wave Principle}. New York: Elliott.
\bibitem{brier1950verification} Brier, G. W. (1950). Verification of forecasts expressed in terms of probability. \emph{Monthly Weather Review}, 78(1), 1--3.
\bibitem{wu2026supercycle} Wu, Y. (2026). Backtesting Specification: The 2025--2026 Semiconductor Storage Supercycle. \emph{Aberdeen Institute of Data Science and AI, South China Normal University}, Tech Rep. SCNU-IDS-2026-08.
\end{thebibliography}

\end{document}
